\begin{tikzpicture}[font=\small, >=Latex]

\tikzset{
  flow/.style={
    -{Latex[length=2.2mm,width=1.5mm]},
    draw=black!80,
    line width=0.85pt
  },
  panel/.style={
    draw=black!75,
    rounded corners=4pt,
    line width=0.8pt
  },
  titlebox/.style={
    draw=black!70,
    rounded corners=3pt,
    fill=blue!5,
    minimum height=0.78cm,
    inner xsep=8pt,
    align=center
  },
  infobox/.style={
    draw=black!70,
    rounded corners=3pt,
    fill=gray!8,
    text width=3.45cm,
    minimum height=1.05cm,
    inner sep=5pt,
    align=center,
    line width=0.75pt
  },
  casebox/.style={
    draw=black!70,
    rounded corners=3pt,
    fill=white,
    text width=4.65cm,
    minimum height=2.35cm,
    inner sep=6pt,
    align=left,
    line width=0.75pt
  },
  goodcase/.style={
    casebox,
    fill=green!10
  },
  badcase/.style={
    casebox,
    fill=red!8
  },
  note/.style={
    draw=black!70,
    rounded corners=4pt,
    fill=gray!8,
    text width=10.9cm,
    minimum height=1.0cm,
    inner sep=6pt,
    align=center,
    line width=0.8pt
  }
}

% ==========================================================
% Left panel
% ==========================================================

\draw[panel, fill=orange!10]
  (0,3.25) rectangle (4.45,10.55);

% \node[
%   titlebox,
%   minimum width=3.3cm
% ] at (2.225,10.95)
% {
%   旧策略与裁剪区间
% };

\node[infobox] at (2.225,9.45)
{
  旧策略\\[0.5mm]
  $\pi_{\theta_{\mathrm{old}}}(B\mid s_t)=0.40$
};

\node[infobox] at (2.225,8.02)
{
  裁剪阈值\\[0.5mm]
  $\epsilon=0.2$
};

\node[infobox] at (2.225,6.55)
{
  概率比值允许范围\\[0.5mm]
  $[1-\epsilon,1+\epsilon]=[0.8,1.2]$
};

\node[
  infobox,
  text width=3.65cm,
  minimum height=1.6cm
] at (2.225,4.78)
{
  裁剪限制的是\\[0.5mm]
  超出该范围后的\\[0.5mm]
  目标函数额外收益
};

% ==========================================================
% Arrows from left panel
% ==========================================================

\draw[flow]
  (4.45,8.55)
  --
  (5.55,8.55);

\draw[flow]
  (4.45,5.25)
  --
  (5.55,5.25);

% ==========================================================
% Positive advantage panel
% ==========================================================

\draw[panel, fill=green!8]
  (5.7,7.15) rectangle (17.0,10.55);

% \node[
%   titlebox,
%   minimum width=3.8cm
% ] at (11.35,10.95)
% {
%   正优势动作 $\widehat{A}_t>0$
% };

\node[casebox] (posleft) at (8.75,8.88)
{
  \textbf{新策略概率：}
  $0.48$\\[0.8mm]

  \textbf{概率比值：}\\
  $r_t(\theta)=0.48/0.40=1.20$\\[0.8mm]

  \textbf{结论：}\\
  位于上界，允许提高动作概率
};

\node[goodcase] (posright) at (14.05,8.88)
{
  \textbf{新策略概率：}
  $0.60$\\[0.8mm]

  \textbf{概率比值：}\\
  $r_t(\theta)=0.60/0.40=1.50$\\[0.8mm]

  \textbf{结论：}\\
  超出上界 $1.20$，裁剪后按 $1.20$ 计算收益
};

% \node[
%   text=green!45!black,
%   font=\footnotesize
% ] at (8.75,7.52)
% {
%   适度提高概率
% };

% \node[
%   text=green!45!black,
%   font=\footnotesize
% ] at (14.05,7.52)
% {
%   过大变化不再获得额外鼓励
% };

% ==========================================================
% Negative advantage panel
% ==========================================================

\draw[panel, fill=red!6]
  (5.7,3.25) rectangle (17.0,6.65);

% \node[
%   titlebox,
%   minimum width=3.8cm
% ] at (11.35,7.05)
% {
%   负优势动作 $\widehat{A}_t<0$
% };

\node[casebox] (negleft) at (8.75,4.98)
{
  \textbf{新策略概率：}
  $0.32$\\[0.8mm]

  \textbf{概率比值：}\\
  $r_t(\theta)=0.32/0.40=0.80$\\[0.8mm]

  \textbf{结论：}\\
  位于下界，允许降低动作概率
};

\node[badcase] (negright) at (14.05,4.98)
{
  \textbf{新策略概率：}
  $0.20$\\[0.8mm]

  \textbf{概率比值：}\\
  $r_t(\theta)=0.20/0.40=0.50$\\[0.8mm]

  \textbf{结论：}\\
  低于下界 $0.80$，裁剪后按 $0.80$ 计算收益
};

% \node[
%   text=red!65!black,
%   font=\footnotesize
% ] at (8.75,3.62)
% {
%   适度降低概率
% };

% \node[
%   text=red!65!black,
%   font=\footnotesize
% ] at (14.05,3.62)
% {
%   过大变化不再获得额外鼓励
% };

% ==========================================================
% Bottom note
% ==========================================================

% \node[note] at (11.35,1.75)
% {
%   PPO 并不直接禁止新策略继续变化，\\
%   而是限制超出裁剪区间之后所能获得的额外优化收益
% };

\end{tikzpicture}